% ── Section 3: Dataset Construction ───────────────────────────────────────────
\section{Dataset Construction}
\label{sec:dataset}

\subsection{Design Principles}

BanglaMix is designed around three principles:
\begin{enumerate}[leftmargin=*, itemsep=2pt]
  \item \textbf{Dialectal authenticity}: audio must contain genuine regional
        dialect, not standard Bengali spoken by someone from a particular region.
  \item \textbf{Natural code-switching}: CS must arise spontaneously in
        informal speech, not from scripted content.
  \item \textbf{Scalability}: the pipeline must be fully automated to allow
        future expansion without human transcribers.
\end{enumerate}

\subsection{Data Sources}

We source audio from two platforms:
\begin{itemize}[leftmargin=*, itemsep=2pt]
  \item \textbf{YouTube} (primary): creator vlogs, technology reviews, comedy
        and roast videos, reaction videos, and campus/career discussions.
  \item \textbf{Facebook} (supplementary): short-form videos and reels,
        particularly for dialects underrepresented on YouTube (Comilla,
        Old~Dhaka, Noakhali).
\end{itemize}

YouTube is downloaded using \texttt{yt-dlp}~\citep{ytdlp} with Bengali-script
search queries targeting dialect-name terms.
Facebook downloads use the same tool with direct URLs and browser cookie
authentication.

\paragraph{Query design.}
A critical design decision is the search query strategy.
Region-name queries (e.g., ``চট্টগ্রামের vlog'', ``Chittagong travel'')
retrieve videos of people \emph{from} that region who predominantly speak
standard Bengali for a general audience.
Instead, we use \emph{dialect-name} queries (e.g., ``চাটগাঁইয়া ভাষায়
comedy english'') that only surface videos where creators consciously use the
dialect---a much stronger signal for genuine dialectal content.
We additionally use the phrase ``আঞ্চলিক ভাষায়'' (in regional dialect) as a
search signal.

For each dialect we design 12 queries across five content angles:
\begin{enumerate}[label=\Alph*., leftmargin=2em, itemsep=1pt]
  \item Dialect-name + technology / phone review (elicits CS with technical
        English terms)
  \item Dialect-name + comedy / roast / prank (natural unscripted dialect;
        comedy cannot sustain a fake accent)
  \item Dialect-name + reaction / gaming (spontaneous commentary)
  \item ``আঞ্চলিক ভাষায়'' + campus / career (CS with academic English)
  \item ``আঞ্চলিক ভাষায়'' + vlog / day-in-my-life (mixed informal CS)
\end{enumerate}

Quality filters \texttt{--dateafter 20210101} and
\texttt{--match-filter "view\_count > 500"} are applied to exclude pre-2021
content (which is less likely to contain CS) and zero-audience uploads.

\subsection{Dialects and Coverage}

Table~\ref{tab:dataset_stats} shows the 15 dialects covered, with clip counts
and durations after segmentation.

\begin{table}[h]
\centering
\caption{BanglaMix dialect statistics after segmentation (Step 2).}
\label{tab:dataset_stats}
\small
\begin{tabular}{llrr}
\toprule
\textbf{Dialect} & \textbf{Region} & \textbf{Clips} & \textbf{Hours} \\
\midrule
Sylhet       & Northeast Bangladesh + UK diaspora & 9{,}665 & 31.85 \\
Old\_Dhaka   & Historic Dhaka city               & 5{,}767 & 17.84 \\
Sandwip      & Island, Chittagong coast           & 5{,}744 & 16.31 \\
Barishal     & South-central Bangladesh           & 4{,}048 & 13.57 \\
Khulna       & Southwest Bangladesh               & 3{,}540 & 11.40 \\
Chittagong   & Southeast Bangladesh               & 3{,}661 & 11.16 \\
Kishoreganj  & Central Bangladesh                 & 3{,}038 &  9.52 \\
Comilla      & Southeast Bangladesh               & 2{,}724 &  9.39 \\
Narail       & Southwest Bangladesh               & 2{,}714 &  8.93 \\
Rangpur      & Northwest Bangladesh               & 2{,}650 &  8.67 \\
Narsingdi    & Central Bangladesh                 & 2{,}150 &  6.63 \\
Mymensingh   & North-central Bangladesh           & 1{,}831 &  5.50 \\
Habiganj     & Northeast Bangladesh               & 1{,}694 &  5.28 \\
Tangail      & Central Bangladesh                 & 1{,}396 &  4.77 \\
Noakhali     & Southeast Bangladesh               & 1{,}235 &  3.93 \\
\midrule
\textbf{Total} & 15 dialects, 9 districts & \textbf{51{,}857} & \textbf{164.74} \\
\bottomrule
\end{tabular}
\end{table}

\subsection{Audio Processing Pipeline}

\paragraph{Step 1 --- Download.}
Audio is extracted at 16~kHz mono WAV using \texttt{yt-dlp} with
\texttt{ffmpeg}. A shared deduplication archive prevents the same video being
downloaded twice across dialect folders.

\paragraph{Step 2 --- Voice Activity Detection and Segmentation.}
Long audio files are segmented into 5--15~second clips using WebRTC~VAD
\citep{webrtcvad} at aggressiveness level~2.
Speech frames are merged with a 300~ms padding on each boundary.
Clips shorter than 5~seconds are discarded; clips longer than 15~seconds are
split into equal-length sub-clips.
This produces 51,857 clips with a mean duration of $11.4 \pm 2.8$~seconds.

\paragraph{Step 3 --- Transcription and CS Detection.}
Each clip is transcribed using Whisper large-v3 \citep{radford2023whisper}
with forced language decoding (\texttt{language=``bn''}).
Three filtering rules are applied:
\begin{enumerate}[leftmargin=2em, itemsep=1pt]
  \item \textbf{Noise filter}: clips where the average \texttt{no\_speech\_prob}
        $\geq 0.60$ are discarded.
  \item \textbf{Language filter}: clips whose transcript contains more
        Hindi/Telugu/Tamil Unicode characters than Bengali characters are
        discarded.
  \item \textbf{CS detection}: each remaining clip is classified using
        word-level script tags. A clip is labelled \textbf{CS} if it contains
        $\geq 2$ Bengali words, $\geq 2$ English words, $\geq 1$ language
        boundary, and an English word ratio of 5--70\%. Otherwise it is
        labelled \textbf{BN\_ONLY}.
\end{enumerate}

CS detection exploits the fact that Bengali and Latin scripts are visually
distinct Unicode ranges ($\backslash$u0980--$\backslash$u09FF vs. ASCII
a--z), making word-level language identification near-perfect without a
dedicated LID model.
We validate the CS labels by manually verifying a random sample of 200 clips
(see Section~\ref{sec:analysis}).

\paragraph{Step 4 --- Silver-standard annotation with LLM ensemble validation.}
We apply a two-stage annotation strategy that eliminates the need for human
transcribers entirely.

\textit{Stage 4a --- Automatic acceptance.}
High-confidence Whisper transcripts (\texttt{no\_speech\_prob} $< 0.30$,
$\geq 3$ Bengali characters in output) are accepted directly as silver-standard
labels~\citep{gandhi2023whisper}.
Clips with \texttt{no\_speech\_prob} $\geq 0.60$ are rejected as noise.

\textit{Stage 4b --- LLM ensemble judgment for borderline clips.}
The remaining clips ($0.30 \leq$ \texttt{no\_speech\_prob} $< 0.60$),
which Whisper is uncertain about, are evaluated by a three-model LLM ensemble
using Mistral-7B \citep{jiang2023mistral}, LLaMA-3.1-8B \citep{dubey2024llama},
and Gemma-2-9B \citep{team2024gemma}, all running locally via Ollama.
Each model independently scores the transcript on four dimensions:
(1) Bengali script validity,
(2) code-switching naturalness,
(3) dialect authenticity, and
(4) absence of noise or hallucination.
A clip is \textbf{accepted} if $\geq 2$ of the 3 models vote ACCEPT,
\textbf{rejected} if $\geq 2$ vote REJECT, and flagged as \textbf{uncertain}
otherwise.
Uncertain clips are excluded from training data (conservative approach).

This ensemble design mitigates individual LLM biases: a single model may
be inconsistent in judging dialectal Bengali, but the majority vote across
three architecturally distinct models is substantially more reliable.
Running models locally preserves data privacy and incurs no API costs.

\paragraph{Step 5 --- Transcript normalisation.}
Before building the final dataset, all transcripts are processed through a
two-level normalisation pipeline implemented in \texttt{text\_normalization.py}:

\begin{enumerate}[leftmargin=2em, itemsep=1pt]
  \item \textbf{Surface cleaning} (applied to both training and evaluation
        transcripts): Unicode NFC normalisation; removal of zero-width
        characters (ZWNJ, BOM), URLs, hashtags, @mentions, and emoji;
        Bengali digit to ASCII digit conversion ($\text{০} \to \text{0}$);
        repeated character collapse (\textit{হাহাহাহা} $\to$ \textit{হাহা});
        punctuation stripping and whitespace normalisation.
  \item \textbf{Pronunciation variant mapping} (applied during evaluation
        only): maps phonetic and script variants of borrowed English words to
        a single canonical form using a curated lexicon of \TODO{200+} entries
        (e.g., \textit{miting} $\to$ \textit{meeting},
        \textit{মিটিং} $\to$ \textit{meeting},
        \textit{ofis} $\to$ \textit{office}).
        This prevents WER inflation due to legitimate spelling variation of
        English loanwords rather than genuine recognition errors.
        Models are trained on the \emph{original} transcript; only evaluation
        references and hypotheses are normalised.
\end{enumerate}

The normalisation-aware WER is reported alongside raw WER in all tables.

\paragraph{Step 6 --- Dataset splits.}
Clips are split at the \emph{source level}: all clips from the same original
video are kept in the same split to prevent speaker leakage across train/dev/test.
An 80/10/10 ratio with dialect-stratified source-level sampling ensures every
dialect appears in all three splits proportionally.

\begin{table}[h]
\centering
\caption{BanglaMix train/dev/test split statistics.}
\label{tab:splits}
\small
\begin{tabular}{lrrrr}
\toprule
\textbf{Split} & \textbf{Clips} & \textbf{Hours} & \textbf{CS clips} & \textbf{CS\%} \\
\midrule
Train & 33{,}568 & 105.15 & 1{,}940 & 5.8\% \\
Dev   & 2{,}768  & 8.89   & 185    & 6.7\% \\
Test  & 5{,}163  & 17.18  & 236    & 4.6\% \\
\midrule
Total & 41{,}499 & 131.23 & 2{,}361 & 5.7\% \\
\bottomrule
\end{tabular}
\end{table}


\subsection{Data Augmentation}

With \TODO{XX,XXX} CS clips in the training set, our corpus is already
substantially larger than prior CS-ASR benchmarks (MUCS~2021: $\sim$5~hours
of CS speech; SEAME: 60~hours).
We therefore do not apply synthetic data generation in our primary experiments.
However, to study robustness to data scarcity we conduct an auxiliary
experiment in which audio augmentation is applied to the CS training subset:
speed perturbation (factors $\times 0.9$ and $\times 1.1$), additive white
noise (20~dB SNR), and pitch shift ($\pm2$ semitones).
Results of this ablation are reported in Appendix~A.
Augmented clips are added to the training set only; the dev and test sets
remain unaugmented.

\subsection{Dataset Quality}

To validate silver-standard transcript quality we manually transcribed a
random sample of 200 CS clips and computed human-verified WER against the
Whisper-generated labels.
Results are reported in Section~\ref{sec:analysis}.
We also verify dialect authenticity by having a native speaker of each dialect
listen to 20 randomly sampled clips and confirm that the clips contain genuine
dialectal speech rather than standard Bengali.
